% how_to_use.tex - Guide to reading and using this book
% ============================================================================

\chapter*{How to Use This Book}
\addcontentsline{toc}{chapter}{How to Use This Book}

\section*{Read in Order (At First)}

The chapters build on each other. Chapter 2 assumes you completed Chapter 1. Skipping ahead usually leads to confusion and backtracking. Once you've worked through the whole book, it becomes a useful reference for jumping to specific topics. But for first-time learners: start at the beginning.

\section*{Special Boxes}

Throughout the book, you'll find colored boxes that highlight different types of information:

\begin{warning}
\textbf{Warnings} alert you to common mistakes and pitfalls. Pay attention to these---they'll save you debugging time.
\end{warning}

\begin{tip}
\textbf{Tips} offer shortcuts, best practices, or alternative approaches that experienced programmers use.
\end{tip}

\begin{note}
\textbf{Notes} provide additional context or clarification that might help your understanding but isn't strictly necessary.
\end{note}

\begin{osbox}
\textbf{OS Differences} explain when something works differently on Windows, macOS, or Linux. Most of Python is identical across systems, but installation and file paths sometimes differ.
\end{osbox}

\begin{labexample}
\textbf{Lab Examples} connect programming concepts to real scientific scenarios. These show how you might actually use what you're learning in your research.
\end{labexample}

\begin{ponder}
\textbf{Things to Ponder} are reflection questions at the end of sections. You don't need to answer them in writing, but thinking about them helps solidify concepts.
\end{ponder}

\section*{Code Examples}

Code appears in monospace font with syntax highlighting:

\begin{lstlisting}
# This is Python code
temperature = 298.15  # Kelvin
print(f"Temperature: {temperature} K")
\end{lstlisting}

When we show interactive Python sessions (the REPL), you'll see \py{>>>} prompts:

\begin{lstlisting}
>>> 2 + 2
4
>>> "hello".upper()
'HELLO'
\end{lstlisting}

Lines starting with \py{>>>} are what you type. Lines without are Python's output.

\section*{Typing vs. Copying}

Type the code examples yourself. Yes, it's slower than copy-pasting. That's the point. Typing forces you to read every character, notice the syntax, and build muscle memory. It also means you'll make typos, see error messages, and learn to fix them---all valuable skills.

\section*{Data Files}

Some examples use data files (CSV files, text files with measurements, etc.). These are available in the \filepath{data/} folder that accompanies this book. When an example references a file like \filepath{titration\_data.csv}, you'll find it there.

\section*{Getting Unstuck}

When something doesn't work:

\begin{enumerate}
    \item Read the error message carefully. Python tells you what's wrong.
    \item Check for typos. Case matters. Indentation matters. Quotes must match.
    \item Re-read the relevant section. Something may have not registered the first time.
    \item Search online. Copy-paste the error message into a search engine. Someone else has had the same problem.
    \item Take a break. Fresh eyes find bugs faster.
\end{enumerate}

Debugging is part of programming. Getting stuck doesn't mean you're bad at this---it means you're learning.
